% TOP_PROBLEM: {"id": 2305067, "title": "Research Problems in Function Theory — Problem 5.67", "category_id": 8, "category": {"id": 8, "name": "analysis", "display_name": "Analysis", "description": "Limits, continuity, calculus, and function theory.", "slug": "analysis", "order_index": 8, "created_at": "2026-07-31T15:26:25.671Z"}, "status": "open", "problem_number": "AMR-022-5067", "set_id": 13}
% FIRST_POSTED: 2026-10-01
% ENTRY_KIND: statement_only
% UnsolvedMath ID 2305067; source catalogue code AMR-022-5067
% !TeX program = lualatex
% Definition and sources only; this file is not an LLM solution attempt.
\documentclass[11pt,a4paper]{article}
\usepackage[margin=25mm]{geometry}
\usepackage{fontspec}
\setmainfont{lmroman10-regular.otf}[BoldFont=lmroman10-bold.otf,
 ItalicFont=lmroman10-italic.otf,BoldItalicFont=lmroman10-bolditalic.otf]
\usepackage{amsmath,amssymb,mathtools}
\usepackage{unicode-math}
\setmathfont{latinmodern-math.otf}
\AtBeginDocument{\let\setminus\smallsetminus}
\usepackage{xurl,hyperref}
\hypersetup{colorlinks=true,urlcolor=blue}
\setcounter{secnumdepth}{0}
\setlength{\parindent}{0pt}
\setlength{\parskip}{5pt}
\setlength{\emergencystretch}{3em}
\begin{document}
\section*{Research Problems in Function Theory — Problem 5.67}\label{2305067}
{\small UnsolvedMath ID 2305067 · AMR-022-5067 · Analysis · AMR Open Problem Lists}
\subsection{Definitions and mathematical statement}
Let $f$ be analytic in $\mathbb D$ and represented there by a
lacunary power series
\[
 f(z)=\sum_{k=0}^{\infty}a_kz^{n_k},\qquad
 1\leq n_0<n_1<\cdots,\qquad n_{k+1}/n_k\geq\lambda>1.
\]
The exponents are integers, the coefficients are complex, and
one fixed gap constant $\lambda$ must work for the entire
series. Define its maximum term by
$m_0(r,f)=\max_{k\geq0}|a_k|r^{n_k}$ for $0<r<1$,
and assume $m_0(r,f)\to\infty$ as $r\uparrow1$.

For angles represented in $[0,2\pi)$, put
\[
 E_f=\left\{\theta:
       \liminf_{r\uparrow1}
       \frac{|f(re^{i\theta})|}{m_0(r,f)}>0\right\}.
\]
Must $E_f$ have one-dimensional Lebesgue measure zero for
every such lacunary series? Equivalently, must almost every
radial direction admit radii tending to the boundary along
which this normalized modulus tends to zero?

The normalization is the largest individual term, not the
maximum modulus of $f$ on a circle or an integral mean.
The limit inferior is taken along all real radii approaching
$1$; the witnessing subsequence may depend on the angle.
The desired assertion imposes no further growth comparison
between maximum terms at different radii and no probabilistic
assumption on the coefficients. Directions at which the
normalized modulus stays bounded away from zero may exist,
but the question is whether they can occupy a set of
positive angular measure.
\subsection{Short English statement}
For a lacunary analytic series with growing maximum term, is its modulus arbitrarily small relative to that term along almost every radius?
\subsection{Sources}
\begin{itemize}
\item[S1] ULAM.AI and the UnsolvedMath Contributors, supplied problems.json, record 2305067 (AMR-022-5067), AMR Open Problem Lists. Catalogue identity and source transcription; CC BY 4.0.
\url{https://huggingface.co/datasets/ulamai/UnsolvedMath}.
\item[S2] Hayman - Research Problems in Function Theory (2018), Problem 5.67. Original source indexed by AMR-022-5067.
\url{https://arxiv.org/abs/1809.07200}.
\end{itemize}
\end{document}
